\section{Comparison with an independently balanced metric}\label{sec:calibration}

The reduced finite equation differs from quadratic balance by the
residual of Proposition~\ref{prop:projection}. Its derivative may have
arbitrarily small positive eigenvalues. We obtain the required metric
correction by excluding zeros on the boundary of a small ball, rather
than estimating the inverse of that derivative.

\begin{theorem}\label{thm:calibration}
There are constants $A,\sigma_0>0$ such that the following holds
uniformly in the fixed background neighborhood. Let $Z$ satisfy
$\mu_s(Z)=b(s)$, let its complex centralizer be scalar, and suppose
$0<\sigma=q_s(Z)<\sigma_0$. Then some $Y\in\hh$ satisfies
\begin{equation}\label{eq:calibration}
\|Y\|_{\op}\le A\sigma,\qquad
\mathcal F(s,e^{Y/2}\cdot Z)=0.
\end{equation}
For $Z'=e^{Y/2}\cdot Z$, the metric $h_{0,s}e^{u(s,Z')}$ is HE
on $E_{s,Z'}$, and $\|u(s,Z')\|_{C^a}\le C_a\sigma$ for every
fixed $a$. The constants do not depend on a positive finite
eigenvalue or on a ratio of such eigenvalues.
\end{theorem}

\begin{proof}
On the real space $\hh$ consider the homotopy into $\mathcal W$
\begin{equation}\label{eq:homotopy}
H_\tau(Y)=\mu_s(e^{Y/2}\cdot Z)-b+
\tau\mathcal R(s,e^{Y/2}\cdot Z),\qquad 0\le\tau\le1.
\end{equation}
Use the ball $\|Y\|_{\op}<a$, where $a=A\sigma<1$ and
$\|Y\|_{\op}=\max_v\|Y_v\|_{\op}$. Write
$M=\sum_vm_v$. Both layers are nonzero, so $M\ge2$.

Suppose $\|Y\|_{\op}=a$. Diagonalize each block of $Y$ unitarily
and sort all $M$ eigenvalues as $y_1\le\cdots\le y_M$, retaining
their vertices. The condition $\sum_vr_v\tr Y_v=0$, with positive
weights, implies that their oscillation is at least $a$. An adjacent
gap therefore satisfies
\begin{equation}\label{eq:gapchoice}
\delta=y_{\ell+1}-y_\ell\ge\frac{a}{M-1}.
\end{equation}
Let $\Pi$ be the tuple of spectral projections onto eigenvectors
above this gap. It is nonzero and proper and commutes with $Y$.

Use orthonormal arrow coordinates at this fixed $s$. If $p$ is a
target eigenvector and $q$ a source eigenvector, write $w_{pq}$ for
the sum of squared arrow entries joining them. The gauge
$e^{Y/2}$ multiplies $w_{pq}$ by $e^{y_p-y_q}$. Since $\mu_s(Z)=b$,
\begin{align}
\langle\Pi,\mu_s(e^{Y/2}\cdot Z)-b\rangle
&=\sum_{p,q}(\Pi_p-\Pi_q)(e^{y_p-y_q}-1)w_{pq}\notag\\
&\ge(1-e^{-\delta})q_s([\Pi,Z])\notag\\
&\ge e^{-2a}\frac{a}{M-1}q_s([\Pi,Z]).\label{eq:boundarypositive}
\end{align}
For a high target and low source the exponential difference is
positive and at least $e^\delta-1$. For a low target and high
source the projection sign is negative and its product with the
exponential difference is at least $1-e^{-\delta}$. These are all
crossing entries; preserving entries contribute zero. The last
inequality follows from $\delta\le2a$ and
$1-e^{-\delta}\ge\delta e^{-\delta}$.
The crossing energy is strictly positive. Otherwise $\Pi$ would
commute with every arrow and be a nonscalar centralizer element.

Put $Z_Y=e^{Y/2}\cdot Z$. The gauge bounds and $[\Pi,Y]=0$ give
\begin{equation}\label{eq:crossingbounds}
q_s(Z_Y)\le e^{2a}\sigma,\qquad
q_s([\Pi,Z_Y])\le e^{2a}q_s([\Pi,Z]).
\end{equation}
Proposition~\ref{prop:projection} consequently yields
\begin{equation}\label{eq:boundaryerror}
|\langle\Pi,\mathcal R(s,Z_Y)\rangle|
\le Ce^{4a}\sigma q_s([\Pi,Z]).
\end{equation}
Choose $A>2C(M-1)$, and then choose $\sigma_0$ so small that
$e^{6A\sigma_0}<2$ and the entire gauge ball stays in the
coefficient neighborhood. Since $a=A\sigma$,
\eqref{eq:boundarypositive} strictly dominates
\eqref{eq:boundaryerror}. Thus $H_\tau$ has no zero on the
boundary, for any $\tau\in[0,1]$. The projection depends on $Y$;
only pointwise exclusion of a boundary zero is needed.

At $\tau=0$ the unique zero is $Y=0$. Indeed the function
\eqref{eq:convexpotential} has $f'(0)=0$ and $f''(t)>0$ for
nonzero $Y\in\hh$, by the scalar-centralizer hypothesis.
It follows that $\langle Y,H_0(Y)\rangle=f'(1)>0$.
The derivative at zero also satisfies
\begin{equation}\label{eq:finitehessian}
\langle Y,DH_0(0)Y\rangle=q_s([Y,Z])>0\quad(Y\ne0).
\end{equation}
Under the fixed identification $T_v\mapsto T_v/r_v$ from
$\mathcal W$ to $\hh$, this is a positive definite real Hessian.
Its local Brouwer degree is $+1$. The operator-norm ball is a
bounded open subset of a finite-dimensional real vector space;
smoothness of its boundary is not required for degree. Homotopy
invariance gives a zero of $H_1$ inside this ball.

At that zero the reduced equation and the complementary equation
both vanish, so \eqref{eq:conjugatedcurvature} is zero. This is
the HE equation. Finally \eqref{eq:crossingbounds} and
Lemma~\ref{lem:complement} give the asserted bounds for $u$.
\end{proof}

The positivity in \eqref{eq:finitehessian} is used to compute a
degree at each fixed parameter. No uniform lower bound for this
Hessian is used. Likewise the crossing energy in
\eqref{eq:boundarypositive} may be arbitrarily small; it appears
as the same positive factor in the error bound and cancels in the
comparison. This gives the uniformity at weak directions.

\begin{proposition}\label{prop:actual-stability}
After shrinking the parameter neighborhood, $E$ is slope-stable
exactly at the $b(s)$-stable points of its representation.
At these points the metric comparison \eqref{eq:main-metric} holds.
\end{proposition}

\begin{proof}
Let $W$ be a proper nonzero subrepresentation. The smooth subbundle
$G_W=\bigoplus_vF_v\otimes W_v$ is preserved by
$\dbar_{s,Z^0}$ and hence is an actual holomorphic subbundle of $E$
with locally free quotient. Put $r_W=\sum_vr_v\dim W_v$.
Its degree is the sum of the degrees of its factor copies, so
\begin{equation}\label{eq:subbundle-slope}
\kappa(G_W)-\kappa_E
=-\frac{\sum_vb_v\dim W_v}{r_W}.
\end{equation}
Stability of $E$ forces the inequalities
\eqref{eq:finite-stability}. This implication uses only the
displayed subbundles, without asserting that they exhaust all
subsheaves.

Conversely, finite stability gives a balanced $Z$ by
Proposition~\ref{prop:balance}. Lemma~\ref{lem:algebra} gives
its scalar centralizer. Since $b(0)=0$, shrink the background
neighborhood so that every such $\sigma$ is below the fixed
threshold of Theorem~\ref{thm:calibration}. That theorem constructs
an HE metric on $E_{s,Z'}$, and the harmonic and finite gauges pull
it back to $E$. An HE bundle is polystable
\cite[Chapter~V]{kobayashi1987}; a simple polystable bundle has
only one stable summand, and therefore is stable. This argument
constructs the HE metric before using stability sufficiency.

For the comparison, work first on the independently balanced
presentation. With $g=e^{Y/2}$ the additional pullback of the
metric has matrix $g^*e^ug$ in a split unitary frame. Scalar
Hermitian-form bounds for $e^u$ and for $g^*g=e^Y$ give
\[
e^{-C\sigma}h_{0,s}\le g^*(h_{0,s}e^u)
\le e^{C\sigma}h_{0,s}.
\]
No commutation of $u$ and $Y$ is needed. Pullback through the
same original balancing gauge and harmonic identification
preserves these inequalities, even if the balancing gauge is
unbounded. This proves \eqref{eq:main-metric}. The unipotent
identification and determinant-normalized finite gauges have
determinant one; $\sum_vr_v\tr Y_v=0$ and $\tr u=0$ preserve
the fixed determinant metric. The HE metric is unique after
this normalization.
\end{proof}

No smooth selection of $Y$ as the parameter approaches a chamber
face is asserted. The proof supplies a uniformly small correction
at every stable point. Bounds for derivatives of a selected gauge
with respect to the parameter are a different question.
